{\color{ink}\bfseries 例 5-6}\par\vspace{1bp}所列操作顺序为{\symfont ②}\(\to \){\symfont ⑥}\(\to \){\symfont ④}\(\to \){\symfont ③}\(\to \){\symfont ①}\(\to \){\symfont ⑤}，故{\symfont ①}前后为{\symfont ③}、{\symfont ⑤}，{\symfont ⑥}后为{\symfont ④}。{\symfont ②}使P阻塞并让出CPU，{\symfont ①}唤醒后P就绪，{\symfont ⑤}须等P重新被调度运行后执行；{\symfont ③}属于键盘驱动，中断处理运行于内核态。\par {\color{ink}图示见合订本第145、146、147页侧栏。}
